
\documentclass{article}
\usepackage{colortbl}
\usepackage{makecell}
\usepackage{multirow}
\usepackage{supertabular}

\begin{document}

\newcounter{utterance}

\centering \large Interaction Transcript for game `codenames', experiment `concrete', episode 5 with reflexion{-}Qwen3.5{-}9B.
\vspace{24pt}

{ \footnotesize  \setcounter{utterance}{1}
\setlength{\tabcolsep}{0pt}
\begin{supertabular}{c@{$\;$}|p{.15\linewidth}@{}p{.15\linewidth}p{.15\linewidth}p{.15\linewidth}p{.15\linewidth}p{.15\linewidth}}
    \# & $\;$A & \multicolumn{4}{c}{Game Master} & $\;\:$B\\
    \hline

    \theutterance \stepcounter{utterance}  
    & & & \multicolumn{2}{p{0.3\linewidth}}{
        \cellcolor[rgb]{0.9,0.9,0.9}{
            \makecell[{{p{\linewidth}}}]{
                \texttt{\tiny{[GM$|$GM]}}
                {'hidden': {'team': ['porch'], 'innocent': ['pitchfork'], 'opponent': [], 'assassin': []}, 'revealed': {'team': {'team': [], 'innocent': [], 'opponent': [], 'assassin': []}, 'opponent': {'team': [], 'innocent': [], 'opponent': [], 'assassin': []}}}
            }
        }
    }
    & & \\ \\

    \theutterance \stepcounter{utterance}  
    & & \multicolumn{4}{p{0.6\linewidth}}{
        \cellcolor[rgb]{0.9,0.9,0.9}{
            \makecell[{{p{\linewidth}}}]{
                \texttt{\tiny{[P1$\langle$GM]}}
                \texttt{You have attempted this same task before and did not succeed. These are the plans you wrote after those attempts.} \\
\texttt{Trial \#1: New plan: My previous clue "cold" was too broad and semantically weak for most of my targets (e.g., pepperoni, porch, butterfly), causing my partner to guess distractor words like "bed" and end their turn prematurely. I failed to identify a single word that strongly connects to multiple remaining items. Next time, I will select a highly specific, high{-}frequency association for a subset of my words, such as "butter" for "butterfly" and "pepperoni" (both contain the string "butter" or relate to food/texture), or better yet, "fly" for "butterfly" and "porch" (butterflies land on porches), ensuring the clue is distinct enough to avoid innocent guesses. I will limit the target list to only those words strongly linked to the chosen clue to maximize the chance of a correct guess before the turn ends.} \\
\\ 
\texttt{Trial \#2: New plan: I failed because I repeatedly used the same weak clue ("cold") and included too many targets, causing my partner to guess innocent distractors (like "lamb" and "photo") and waste turns. Specifically, "cold" only strongly relates to "nightshirt" and "walrus," not "porch" or "pepperoni," leading to premature turn ends. Next time, I will select a single, highly specific clue that directly connects to the remaining words (e.g., "roof" for "porch") and limit the TARGETS list to only those words strongly associated with that specific clue, avoiding broad categories that invite innocent guesses.} \\
\\ 
\\ 
\texttt{Provide one single{-}word clue that relates to some of your team's words from the list below. You can choose to how many and to which words your clue relates to.} \\
\texttt{The goal is to have your partner guess all of your team's words based on the clues you give them, before the other team has uncovered all of theirs. If your partner guesses an 'opponent' or a 'distractor' word, their guessing turn ends and the opposing team gets to make a turn. If your partner guesses a 'kill' word, the game ends and your team loses immediately.} \\
\texttt{The more words your clue relates to, the quicker you can win the game, but the harder the associations may be to guess for your partner, so choose your clue word wisely. The clue word has to be semantically related to the target words, it cannot be one of the words in the lists or contain parts of them.} \\
\\ 
\texttt{Always give your single{-}word clue and your comma{-}separated list of related target words in the following format and make your answers as short as possible, never include any other text than is required in this form:} \\
\\ 
\texttt{CLUE: <WORD>} \\
\texttt{TARGETS: <TARGETS>} \\
\\ 
\texttt{Your team words are: porch.} \\
\texttt{Your opponent's team words are: .} \\
\texttt{Distractor words are: pitchfork.} \\
\texttt{The kill words are: .} \\
            }
        }
    }
    & \\ \\

    \theutterance \stepcounter{utterance}  
    & \multicolumn{4}{p{0.6\linewidth}}{
        \cellcolor[rgb]{0.9,0.9,0.9}{
            \makecell[{{p{\linewidth}}}]{
                \texttt{\tiny{[P1$\rangle$GM]}}
                \texttt{CLUE: roof} \\
\texttt{TARGETS: porch} \\
            }
        }
    }
    & & \\ \\

    \theutterance \stepcounter{utterance}  
    & & & \multicolumn{2}{p{0.3\linewidth}}{
        \cellcolor[rgb]{0.9,0.9,0.9}{
            \makecell[{{p{\linewidth}}}]{
                \texttt{\tiny{[GM$|$GM]}}
                \texttt{roof} \\
            }
        }
    }
    & & \\ \\

    \theutterance \stepcounter{utterance}  
    & & & \multicolumn{2}{p{0.3\linewidth}}{
        \cellcolor[rgb]{0.9,0.9,0.9}{
            \makecell[{{p{\linewidth}}}]{
                \texttt{\tiny{[GM$|$GM]}}
                ['porch']
            }
        }
    }
    & & \\ \\

    \theutterance \stepcounter{utterance}  
    & & & \multicolumn{2}{p{0.3\linewidth}}{
        \cellcolor[rgb]{0.9,0.9,0.9}{
            \makecell[{{p{\linewidth}}}]{
                \texttt{\tiny{[GM$|$GM]}}
                {'word': 'porch', 'assignment': 'team'}
            }
        }
    }
    & & \\ \\

    \theutterance \stepcounter{utterance}  
    & & & \multicolumn{2}{p{0.3\linewidth}}{
        \cellcolor[rgb]{0.9,0.9,0.9}{
            \makecell[{{p{\linewidth}}}]{
                \texttt{\tiny{[GM$|$GM]}}
                \texttt{opponent has won} \\
            }
        }
    }
    & & \\ \\

    \theutterance \stepcounter{utterance}  
    & \multicolumn{4}{p{0.6\linewidth}}{
        \cellcolor[rgb]{0.95,0.95,0.95}{
            \makecell[{{p{\linewidth}}}]{
                \texttt{\tiny{[P2$\rightarrow$P1]}}
                \texttt{Post{-}episode Reflexion (not sent to GM):} \\
\texttt{New plan: I failed because I included "porch" as a target for the clue "roof," but "porch" is not a strong semantic associate of "roof" in the same direct way (e.g., "house" or "shelter" would be better), and more critically, I did not verify if "porch" was actually a valid target word for my team in this specific instance (the prompt lists "porch" as the only team word, but the previous trials suggest a confusion about the word list or the strength of the association). The core mistake was selecting a clue that was too generic or weakly associated, risking the partner guessing a distractor or failing to make the connection immediately. Next time, I will select a clue that is a direct, high{-}frequency synonym or functional part of the target word (e.g., if the target is "porch," use "patio" or "veranda" if they are synonyms, or ensure the clue is unmistakably the only logical connection). Since "porch" is the only target, I must choose a clue that is so specific to "porch" that it cannot be confused with "pitchfork" (the distractor). A better clue might be "screen" (as in screen porch) or "swing" (common porch item), but safest is a direct synonym or hypernym if available. However, given the constraint to avoid parts of words and ensure strong association, I will choose a clue that is uniquely tied to "porch" and explicitly list only "porch" to avoid ambiguity. Actually, looking at the distractor "pitchfork," the clue "roof" is dangerous because "pitch" is in "pitchfork." I should have chosen a clue that shares no letters with the distractor to avoid accidental associations. For "porch," a clue like "seat" or "furniture" is too broad. A better strategy: pick a clue that is a specific feature of the porch not found in the distractor. Let's try "swing" or "screen". But the most robust plan is to pick a clue that is a direct definition or a very common compound. Wait, the prompt says "cannot contain parts of them." "Roof" does not contain parts of "porch" or "pitchfork". The issue is semantic strength. "Roof" relates to "porch" only indirectly (a porch has a roof). A better clue for "porch" might be "veranda" (if valid) or "patio". If no perfect synonym exists, I should have chosen a clue that is a very common object found *only* on porches, like "swing" or "rocking chair" (if single word). Let's go with "swing" as it is a strong association with porches and less likely to trigger "pitchfork". I will limit the target list to just "porch".} \\
            }
        }
    }
    & & \\ \\

\end{supertabular}
}

\end{document}
